\documentclass[12pt, a4paper]{article}
\usepackage{graphicx} % Required for inserting images
\graphicspath{{images/}{images-desafios/}}
\usepackage[a4paper, total={7.5in, 10in}]{geometry}
\usepackage{float} % pra conseguir usar o \begin{figure}[H]
\usepackage{indentfirst} % pra fazer tab n 1 paragrafos dos capitulos
% package pra tabela
\usepackage{multirow}
\usepackage{array}

%pacote pra codigo
\usepackage{listings} 
\usepackage{xcolor}
\definecolor{codegreen}{rgb}{0,0.6,0}
\definecolor{codegray}{rgb}{0.5,0.5,0.5}
\definecolor{codepurple}{rgb}{0.58,0,0.82}
\definecolor{backcolour}{rgb}{0.95,0.95,0.92}

\lstdefinestyle{myCodeStyle}{
	backgroundcolor=\color{backcolour},   
	commentstyle=\color{codegreen},
	keywordstyle=\color{cyan},
	numberstyle=\tiny\color{codegray},
	stringstyle=\color{codepurple},
	basicstyle=\ttfamily\footnotesize,
	breakatwhitespace=false,         
	breaklines=true,                 
	captionpos=b,                    
	keepspaces=true,                 
	%numbers=left,                    
	numbersep=5pt,                  
	showspaces=false,                
	showstringspaces=false,
	showtabs=false,                  
	tabsize=2
}
\lstset{style=myCodeStyle}
\renewcommand{\lstlistingname}{Código}%Muda o termo "Listing":  Listing -> Código
\renewcommand{\lstlistlistingname}{Lista de \lstlistingname s}% List of Listings -> Lista de Códigos


\newcommand{\unit}[1]{\ensuremath{\, \mathrm{#1}}} % comando pra tirar do math mode a unidade de medida

\def\tamanhoGrafico{0.8}

\newcommand{\ohm}{\unit{\Omega} }
\newcommand{\kohm}{\unit{k\Omega}}

\title{
	Laboratório Digital A - PCS3335
	\\ Planejamento da Experiência 7
}

\author{
	Alunos:
	\\ Rafael Hipólito Rodrigues - N.USP: 14604308
	\\ {{PERSON-2}} - N.USP: {{PERSON-ID-2}}
	\\ \\
	Turma 3 - Bancada A6
	\\ Professor: Bruno de Carvalho Albertini
}

%\date{}

\begin{document}
	
	\maketitle
	
	\section*{Planejamento} 
	Adição do Entidade \texttt{Receiver Buffer Register} ao circuito da UART simplificada.
	Foi mantido as entidades da experiência 6 e 5.
	
	\section*{Diagramas}
	
	\subsection*{FSM}
	
	Diagrama da Maquina de Estados Finitos do Receiver gerado no Quartus : 
	
	\begin{figure}[H]
		\centering
		\includegraphics[width=0.9\linewidth]{images/FSM.png}
		\caption{FSM do Receiver.}
		\label{fig:FSM_exp4}
	\end{figure}
	
	
	
	\subsection*{RTL}
	
	Diagramas RTL da entidades gerado no Quartus:
	
	\begin{figure}[H]
		\centering
		\includegraphics[width=0.9\linewidth]{images/RTL_uart.png}
		\caption{Diagrama RTL da UART.}
		\label{fig:RTL1}
	\end{figure}
	
	\begin{figure}[H]
		\centering
		\includegraphics[width=0.9\linewidth]{images/RTL_receiver.png}
		\caption{Diagrama RTL do \texttt{Receiver}.}
		\label{fig:RTL_TCC}
	\end{figure}

	
	Os contadores e registradores utilizados foram os mesmos da experiência 2.
	
	\section*{Pinagem}
	
	\begin{table}[H]
		\begin{center}
			\begin{tabular}{|l|l|l|}
				\hline
				clock50M           & PIN\_M9   & Input  \\ \hline
				display7seg{[}0{]} & PIN\_AA22 & Output \\ \hline
				display7seg{[}1{]} & PIN\_Y21  & Output \\ \hline
				display7seg{[}2{]} & PIN\_Y22  & Output \\ \hline
				display7seg{[}3{]} & PIN\_W21  & Output \\ \hline
				display7seg{[}4{]} & PIN\_W22  & Output \\ \hline
				display7seg{[}5{]} & PIN\_V21  & Output \\ \hline
				display7seg{[}6{]} & PIN\_U21  & Output \\ \hline
				leds{[}0{]}        & PIN\_AA2  & Output \\ \hline
				leds{[}1{]}        & PIN\_AA1  & Output \\ \hline
				leds{[}2{]}        & PIN\_W2   & Output \\ \hline
				leds{[}3{]}        & PIN\_Y3   & Output \\ \hline
				leds{[}4{]}        & PIN\_N2   & Output \\ \hline
				leds{[}5{]}        & PIN\_N1   & Output \\ \hline
				leds{[}6{]}        & PIN\_U2   & Output \\ \hline
				leds{[}7{]}        & PIN\_U1   & Output \\ \hline
				leds{[}8{]}        & PIN\_L2   & Output \\ \hline
				leds{[}9{]}        & PIN\_L1   & Output \\ \hline
				rbrRead            & PIN\_M9   & Input  \\ \hline
				reset              & PIN\_B16  & Input  \\ \hline
				serialIn           & PIN\_M16  & Input  \\ \hline
				serialOut          & PIN\_N16  & Output \\ \hline
				switches{[}0{]}    & PIN\_U13  & Input  \\ \hline
				switches{[}1{]}    & PIN\_V13  & Input  \\ \hline
				switches{[}2{]}    & PIN\_T13  & Input  \\ \hline
				switches{[}3{]}    & PIN\_T12  & Input  \\ \hline
				switches{[}4{]}    & PIN\_AA15 & Input  \\ \hline
				switches{[}5{]}    & PIN\_AB15 & Input  \\ \hline
				switches{[}6{]}    & PIN\_AA14 & Input  \\ \hline
				switches{[}7{]}    & PIN\_AA13 & Input  \\ \hline
				switches{[}8{]}    & PIN\_AB13 & Input  \\ \hline
				switches{[}9{]}    & PIN\_AB12 & Input  \\ \hline
			\end{tabular}
		\end{center}
	\end{table}
	
	
	
	\section*{Testes}
	
	\subsection*{Resultados dos testes}

	\begin{itemize}
		\item bits adequadamente recebidos em 9600bps.
		\par Resultado: 
		
		\item Ao ativar e desativa o reset, volta ao estado inicial. 
		\par Resultado: 
		
		\item Funcionamento correto do RBR.
		\par Resultado: 
		
		\item Funcionamento dos bits LSR:
		
		\subitem Bit 0: Data Ready (DR) - Resultado: 
		
		\subitem Bit 1: Overrun Error (OE) - Resultado: 
		
		\subitem Bit 2: Parity Error (PE) - Resultado: 
		
		\subitem Bit 3: Framming Error (FE) - Resultado:
		
		\item Valor recebido adequadamente exibido no display de 7 segmentos.
		\par Resultado: 
		
		
		
		
	\end{itemize}
	
	
\end{document}
